\documentclass{article}
\usepackage[utf8]{inputenc}
\usepackage[english]{babel}
\usepackage{amsmath}
\usepackage{amssymb}
\usepackage{geometry}

\geometry{a4paper, margin=1in}

\title{Holistic Logic Theory: Archive of Symbolic and Quantitative Propositions \\ \large (Independent Core System Alphabet)}
\author{Berat Temel}
\date{September 2026}

\begin{document}

\maketitle

\section{Glossary of Variables and Universal Symbols}
The formal definitions of all variables and set operators utilized within the system are structured as follows:
\begin{itemize}
    \item $L$ or $L_s$: \textbf{Set of Human Limitations} (Cognitive capacity and the boundaries of consciousness).
    \item $\mathcal{U}$: \textbf{Categorization Mechanism} (The bounded cognitive universe of human perception).
    \item $q$: \textbf{Human Perception} (Instinctual and cognitive processing capacity).
    \item $E$: \textbf{Absolute Cosmic Universal Set} (The objective reality space independent of human existence).
    \item $R$: \textbf{Observable Reference Point} (Verifiable data belonging to the empirical world).
    \item $\text{val}(x)$: The \textbf{Truth Value} assigned to any given proposition or function.
    \item $u$: \textbf{\textit{Undefined}} / Unknowable truth value (Non-classical logical dimension).
    \item $|p|$: The \textbf{Absolute Reality Value} of a proposition (An algebraic absolute value abstraction).
\end{itemize}

\section{Core Symbolic Propositions and Laws}

\subsection{Proposition 1: The Perceptual Boundary and Categorization Law}
\[ ((p \in \mathcal{U}) \rightarrow q) \land (q \notin E) \]
\textbf{Analysis:} If a phenomenon is categorized by the human mind ($p \in \mathcal{U}$), it is fundamentally bound within the limits of human perception ($q$). Furthermore, this human perception ($q$) cannot be an extension or a part of the absolute cosmic reality, namely the universal set ($E$); it is strictly decoupled from it.

\subsection{Proposition 1.1: The Law of Limit Denial and Systemic Collapse}
\[ (p \rightarrow \neg L_s) \rightarrow \text{val}(p) = 0 \]
\textbf{Analysis:} If the premise of a proposition inherently attempts to deny or transcend the existence of human limitations ($p \rightarrow \neg L_s$), its truth value within our logical space is directly equivalent to zero, meaning it is false ($\text{val}(p) = 0$). Claims that operate outside the boundary of limits are automatically rejected by the system architecture.

\subsection{Proposition 1.2: The Law of Impossibility of Absoluteness Within Boundaries}
\[ (p \in L_s) \rightarrow \neg |p| \]
\textbf{Analysis:} If a proposition is an element of our set of limitations ($p \in L_s$), it is mathematically impossible for that proposition to represent an absolute reality ($\neg |p|$). Any data exposed to a limit loses its absoluteness, collapsing into a mere reflection.

\subsection{Proposition 1.3: Wittgenstein's Fallacy of Independence (The Axiom of Derivation)}
\[ \neg p \rightarrow \neg q \]
\textbf{Analysis:} Contrary to Wittgenstein's assertions, empirical facts are not logically independent of one another. If a core foundational fact is false ($\neg p$), all sub-facts and equations derived from and dependent upon it are automatically condemned to falsehood ($\neg q$). (e.g., If gravity is non-existent, the orbital trajectory equations inevitably collapse).

\subsection{Proposition 1.4: The Reference Deficit Problem and Teistic Collapse}
\[ (g \in \mathcal{T}) \rightarrow \text{val}(T_r) = 0 \quad \text{because} \quad \text{val}(R) = u \]
\textbf{Analysis:} When an transcendent attribute like infinite power ($g$) is integrated into the set of divine characteristics ($\mathcal{T}$), it yields absolutely no corresponding observable reference point ($R$) within the empirical world; thus, its reference value is undefined ($\text{val}(R) = u$). In this system architecture, if the reference of a transcendent existence proposition ($T_r$) is undefined, its truth value is instantly nullified by the system ($\text{val}(T_r) = 0$).

\subsection{Proposition 1.5: Cognitive Determinism and the Illusion of Free Will}
The strict deterministic decision mechanism is expressed through the following functional framework, mapping data captured from the observable space into the personality matrix, thereby neutralizing free will ($D_e$):

\[ K_u(p) \implies \Big( (G_{\ddot{o}} \in G_k) \rightarrow \text{input}(G_{\ddot{o}} \in G_k) \in K_i \Big) \]

The inescapable logical consequence of this systemic input-output relationship is formalized as follows:
\[ K_u \rightarrow \neg D_e \]

\textbf{Variables and Functions:}
\begin{itemize}
    \item $G_k$: The Set of Observability (The baseline database of the external world).
    \item $G_{\ddot{o}}$: The Observed / Decision-Making Agent (Infinitely variable and undefined in quantity).
    \item $K_i$: The Set of Personalities (The internal matrix where the brain logs captured external inputs).
    \item $K_u(p)$: The Decision-Making Process Function (The final dependent outcome variable governed by external inputs).
    \item $D_e$: The Free Will Factor.
    \item $\neg$: The Negation / Olformatting operator.
\end{itemize}

\textbf{Analysis:} The decision-making process ($K_u$) of a human or any deciding agent ($G_{\ddot{o}}$) is not the product of a free, autonomous will. Every single input processed from the external space (The Set of Observability $G_k$) is inevitably written into the Set of Personalities ($K_i$) as a computational input. The decision function ($K_u$) operates entirely as a dependent output driven by these pre-existing inputs. Consequently, the execution of a systemic decision mechanism absolutely confirms the non-existence of free will ($K_u \rightarrow \neg D_e$). The terrain and the observed factors are the true deciders; the human agent is merely the biological hardware running the algorithm.

\section{Digital Electronics and Logical Algebraic Equation}
\[ \text{val}(p) \times \text{val}(q) = \text{val}(p \land q) \]
\textbf{Explanation:} This serves as the mathematical proof demonstrating that the logical "AND" ($\land$) operator functions identically to algebraic multiplication ($\times$), where a single 0 (False) value entering the circuit completely zeroes out the entire system state.

\section{Socio-Economic Wage Functions}

\subsection{General General Labor Wage Formula}
\[ \text{Wage} = \frac{\text{Competence} \times \text{Difficulty} \times \text{Danger}}{3} \]

\subsection{Artist Wage Formula (Artificial Intelligence Shield)}
\[ \text{Wage (Art)} = \frac{\text{Competence} \times \text{Difficulty} \times \text{Danger} \times \text{AI Factor}}{4} \]

\(\subsection{Proposition 1.7: The Trans-Limit (\(L_s\)) and the Totality of Infinity}\)
If a proposition claims that the trans-limit domain can be strictly defined or captured, it triggers a systemic contradiction, proven by the mathematical necessity of an undefinable infinity.

\[\Big( p \rightarrow \Box \Diamond L_s \Big) \rightarrow \neg \Big( p \rightarrow \Box \Diamond L_s \Big) \quad \because \quad \Big( \Diamond \big( \text{val}(p) \in \mathcal{P}_{set} \big) = 0 \Big)\]

\(\textbf{Variables and Modal Operators:} \begin{itemize}     \item\) \(L_s\): \(\textbf{Trans-Limit Space} \)(The reality that exists beyond human cognitive boundaries).
    \item \(\mathcal{P}_{set}\): The Set of Formal Propositions.
    \item \(\Box\): Modal operator for \(\textbf{Necessity} \)(Zorunluluk).
    \item \(\Diamond\): Modal operator for \(\textbf{Possibility} \)(Olasılık).
    \item \(\because\): The logical symbol for \(\textbf{Because} (\)Ters üçgen şeklinde ``çünkü'' operatörü).
    \item $\text{val}(p) = 0$: Systemic invalidation / False state.
\end{itemize}

\textbf{Analysis:} For thousands of years, mathematicians have defined infinity merely as ``that which has no end,'' showcasing the ultimate ceiling of human comprehension. This proposition mathematically proves that if a human premise (\(p\)) asserts a non-hypothetical, fully computed mastery over the post-limit zone (\(\Box \Diamond L_s\)), it instantly negates itself. The mathematical existence of absolute infinity guarantees that there are values we cannot compute. Our cognitive structure possesses an built-in threshold that bars us from processing infinite data. Therefore, the very fact that we encounter an untamable infinity is the absolute, concrete proof of the Limit itself. The limit is verified by our mathematical inability to transcend it.

\subsection{Proposition 1.8: The Epistemic Co-Dependence of Probability and Limitation}
\textbf{Ham Okul Notu (September 2026):} 
Limit nedeniyle olasılıklar doğar. Eğer birçok olasılık varsa, limitte vardır çünkü limitin oluşu nelerin gerçek olup olmadığını belirsiz kılar. Eğer limit var ise olasılık da vardır. Limit yoksa olasılık da yok anlamına gelir ve bu kesin mantık içerisinde tutarsızdır.
q = limit, o = olasılık.

\textbf{Kipsel Mantık Formülü (Modal Axiom of Probability):}
\[ (p \rightarrow q) \land (\Diamond o \rightarrow \Box\Diamond q) \land (\neg q \rightarrow \neg\Diamond o) \]

\textbf{Çözüm Analizi (Epistemic Horizon Proof):}
If the human agent possessed no cognitive threshold (limit), it would comprehend absolute infinity completely. In a state of infinite comprehension, all random variances dissolve into absolute certainty; thus, probability cannot exist for an unbounded mind. Therefore, the very mathematical manifestation of probabilistic outcomes (e.g., Quantum Superposition) is the direct, irreducible proof of the cognitive Limit. Probability is not an ontological randomness, but an epistemic byproduct of our bounded processing capacity.


\end{document}

